
Reducing energy consumption has become a key priority for policymakers seeking to address the climate crisis \citep{shuklaClimateChange20222022}. At the same time, public concern for environmental issues has increased significantly in recent years \citep{andreGloballyRepresentativeEvidence2024,milfontTenyearPanelData2021}. Yet despite increased awareness and widespread intentions to act sustainably, a persistent attitude-behavior gap remains, with households often failing to translate their environmental aspirations into tangible reductions in energy consumption \citep{allcottRethinkingRealtimeElectricity2011,zhangExaminingAttitudebehaviorGap2021}.

A central explanation for this gap is psychological, particularly the salience bias \citep{chettySalienceTaxationTheory2009}. Electricity is typically a low-salience good, as households frequently lack awareness of when and by which appliances electricity is consumed or whether their consumption is particularly intense or concentrated during certain periods \citep{alcocerSalienceBiasFramework2024,jessoeKnowledgeLessPower2014,tiefenbeckRealtimeFeedbackPromotes2019}. While energy-consuming activities, such as heating and lighting, yield immediate and tangible benefits, the associated economic and environmental costs remain abstract and delayed. This discrepancy reinforces a bias toward immediate rewards \citep{kahnemanJudgmentUncertaintyHeuristics1982}, resulting in persistently suboptimal consumption decisions. 

Policymakers increasingly rely on behavioral interventions, particularly personalized feedback, to close the gap between pro-environmental intentions of households and actual energy use. These interventions aim to correct common decision-making biases, such as inattention or salience bias, and can deliver a behavioral double dividend by reducing electricity consumption and improving consumer welfare \citep{allcottShortrunLongrunEffects2014}. Among such tools, energy feedback has received sustained attention for its ability to influence day-to-day usage decisions \citep{agarwalReviewResidentialEnergy2023,khannaMulticountryMetaanalysisRole2021}. The underlying assumption is straightforward: households are more likely to conserve energy when they receive timely and relevant information about their energy use and costs \citep{jessoeKnowledgeLessPower2014}. Feedback mechanisms make the marginal cost of consumption more salient, thereby helping households link specific actions to cost consequences. This can counteract bounded rationality and overconsumption \citep{alcocerSalienceBiasFramework2024}.
In a wide range of studies, feedback has been found to reduce electricity consumption by approximately 2\% to 5\% on average \citep{delmasInformationStrategiesEnergy2013,khannaMulticountryMetaanalysisRole2021,mckerracherEnergyConsumptionFeedback2013}, although the effect sizes vary with the details of the implementation and the levels of engagement.

Despite the growing literature on behavioral energy interventions, most existing evidence stems from short-run field trials, typically lasting only a few weeks or months.\footnote{\cite{delmasInformationStrategiesEnergy2013} reviewed  156 studies and found that close to 60\% of the field studies lasted for three months.} A more recent and comprehensive meta-analysis by \citet{khannaMulticountryMetaanalysisRole2021} covering 122 experimental studies confirms this pattern. The mean treatment duration was just 21.5 weeks, with a median of only 12 weeks. As emphasized by \citet[p.~928]{khannaMulticountryMetaanalysisRole2021}, the scarcity of long-term interventions significantly limits our ability to assess sustained behavioral change and the durability of treatment effects.

This paper addresses these limitations using a randomized controlled trial with 907 households. Participants were randomly assigned to an information condition ($N=454$), which received access to a mobile app that offered  feedback, social comparisons, and gamified elements, or to a control condition ($N=452$) without feedback. Electricity consumption was tracked for up to 70 weeks during treatment and 3 weeks of pre-treatment data. The combination of feedback, social comparison, and gamification was intentionally designed to reduce attrition, a common issue in earlier feedback studies that often suffered from declining user engagement over time \citep{buchananQuestionEnergyReduction2015, fischerFeedbackHouseholdElectricity2008}. Moreover, research suggests that bundling multiple behavioral mechanisms can generate synergistic effects, leading to larger and more persistent changes in consumption behavior than isolated interventions \citep{khannaMulticountryMetaanalysisRole2021}. This design allows us to answer two key questions: (1) Do electricity savings persist over more than a year when households receive ongoing smartphone-based feedback? (2) How does variation in app engagement shape the magnitude and persistence of these energy savings over time? 


The contribution of this paper is twofold. First, it provides rare evidence on behavioral persistence in a high-frequency setting over more than a year—an open gap identified in recent meta-analyses \citep{khannaMulticountryMetaanalysisRole2021}. Second, by leveraging detailed app interaction data, we disentangle average treatment effects from engagement-driven heterogeneity, offering insights into the mechanisms underpinning sustained behavior change. 

Behavioral energy conservation programs have traditionally relied on low-frequency interventions such as Home Energy Reports (HERs), which provide households with comparative usage feedback on a monthly or quarterly basis \citep{allcottSocialNormsEnergy2011}. While HERs have consistently produced statistically significant reductions in electricity consumption, typically in the range of 1\% to 3\%, their behavioral mechanisms are limited by the infrequency and delayed nature of the feedback \citep{allcottShortrunLongrunEffects2014, costaEnergyConservationNudges2013}.

The introduction of in-home displays (IHDs) offered a partial improvement by providing near real-time information directly within the household. Research suggests that IHDs can lead to larger average savings, ranging from 5\% to 10\% \citep{fischerFeedbackHouseholdElectricity2008,khannaMulticountryMetaanalysisRole2021}. However, \citet{khannaMulticountryMetaanalysisRole2021} find that these higher estimates often come from studies with weaker methodological designs, such as small sample sizes, limited external validity, or lack of randomization. A persistent challenge in digital behavioral interventions is user attrition, which can attenuate treatment effects and bias impact estimates upward when analyses fail to account for differential engagement \citep{delmasInformationStrategiesEnergy2013}. To mitigate drop-out and sustain user attention, the app was deliberately designed to combine multiple behavioral levers—including feedback, social comparisons, and gamification. These components have been shown to produce complementary effects \citep{khannaMulticountryMetaanalysisRole2021} and reduce behavioral fatigue over time. 

Our results show a statistically significant intention-to-treat (ITT) effect, with households given access to the app reducing their electricity consumption by approximately 6\% (p $<$ 0.05). This finding significantly exceeds the modest average short-term reductions of around 2\% typically reported in previous studies using in-home displays (IHDs) or online portals \citep{buchananQuestionEnergyReduction2015}, and falls within the upper range of interventions with real-time, personalized feedback, as summarized in \citet{khannaMulticountryMetaanalysisRole2021}. Furthermore, unlike previous interventions, which often experience an action-and-backsliding phenomenon \citep{allcottShortrunLongrunEffects2014,buchananQuestionEnergyReduction2015}. Our dynamic analysis suggests a sustained reduction in electricity use. Specifically, consumption decreased steadily after the initial adjustments, peaked at approximately 15-20 weeks during treatment, and maintained significant reductions for over 30 weeks.  Beyond this period, the information condition continues to consume less electricity on average, but the differences are no longer statistically significant.

This intervention leverages existing smart meter infrastructure and delivers personalized feedback via a mobile app. It demonstrates how scalable, low-cost tools can produce modest, albeit persistent, reductions in electricity consumption. Although the average household savings are small in absolute terms, they add up significantly at scale and are achieved at a much lower cost than hardware-intensive alternatives, such as in-home displays. This makes app-based feedback a highly relevant policy tool for utilities seeking demand-side flexibility and emissions reductions.

In addition, we contribute to the literature by examining variation in treatment effects according to user engagement levels.  Households are categorized as low-, medium-, or high-engagement based on their cumulative number of app interactions during the first 40 weeks of the study.

High-engagement users exhibited the most pronounced and sustained energy savings, maintaining significant reductions throughout the observation period. In contrast, the medium- and low-engagement groups experienced only temporary reductions, returning toward baseline consumption levels within several weeks. Given the overall decline in app usage over time, we acknowledge that sustained reductions among high-engagement users may reflect a threshold effect--where a minimum amount of interaction is necessary to facilitate meaningful learning and lasting behavioral change. In contrast, low levels of initial engagement do not appear sufficient to induce long-term savings. 


The remainder of this paper is organized as follows: Section~\ref{sec:Conceptual_framework} outlines the conceptual framework underlying the study. Section~\ref{sec:Experimental_design} provides a detailed description of the experimental design, data collection, and methodology. Section~\ref{section:Results} presents the results, followed by a discussion of policy implications in Section~\ref{section:Conclusion}.  
